\documentclass[12pt]{article}
\usepackage[utf8]{inputenc}
\usepackage[T1]{fontenc}
\usepackage[margin=1in]{geometry}
\usepackage{setspace}
\usepackage{xcolor}
\usepackage[colorlinks=true, urlcolor=blue!60!black, linkcolor=blue!60!black]{hyperref}
\setlength{\parindent}{0em}
\setlength{\parskip}{0.6em}
\renewcommand{\baselinestretch}{1.0}
\begin{document}
\begin{center}
{\LARGE \textit{McCulloch v. Maryland}, 17 U.S. (4 Wheat.) 316 (1819)\par}
\vspace{0.5em}
{\large Case Brief --- Two-Page Edition\par}
\end{center}

\section{Memory Jogger}
Congress has implied power to charter a national bank, and states may not tax federal instrumentalities (\href{https://supreme.justia.com/cases/federal/us/17/316/}{17 U.S. 316}).

\section{Detailed Case Facts}
Congress chartered the Second Bank of the United States in 1816, and its Baltimore branch opened without Maryland's authorization. Maryland taxed all banks not chartered by its legislature, which in practice meant only the federal branch, requiring stamped paper for notes or an annual \$15,000 payment. Cashier James McCulloch issued unstamped notes and was sued for the penalties.

\section{Procedural History}
The Baltimore County Court ruled against McCulloch and the Maryland Court of Appeals affirmed; the Supreme Court reversed.

\section{Judicial Votes}
Unanimous. Marshall wrote for the Court; no concurrences or dissents (Justice Todd's participation is disputed).

\section{Holding}
Congress may choose any appropriate, unprohibited means plainly adapted to its enumerated powers; ``necessary'' means useful or conducive, not indispensable. States may not tax or otherwise burden the operations of federal instrumentalities.

\section{Analysis of Opinions}
The Constitution derives from the people, not the states, and the Tenth Amendment omits ``expressly.'' A constitution marks only great outlines, and incorporation is merely a means. ``Absolutely necessary'' in Article I, Section 10 shows that ``necessary'' alone is looser; a strict reading would be unworkable; and the Clause sits among Congress's powers, not its limits. A state tax on a federal instrument lets a part control the whole without representation, so a bright-line rule forbids it.

\section{Examples: Future Applications}
\textbf{Same-Side:} federal civil commitment of dangerous prisoners (\textit{Comstock}) is a valid means. \textbf{Opposite-Side:} a general, nondiscriminatory property tax on a federal bank's building is permitted. \textbf{Fence-Sitter:} a felony for falsifying any record of a remote federal subgrantee.

\section{Critique}
Spencer Roane and modern originalists such as Lawson, Natelson, and Baude argue that Marshall loosened ``necessary and proper'' too far; Amar and Balkin defend his reading, and Mikhail and Schwartz think he conceded too much to enumerationism.

\section{Key Quotations}
\begin{quote}
``Let the end be legitimate, let it be within the scope of the Constitution, and all means which are appropriate, which are plainly adapted to that end, which are not prohibited, but consist with the letter and spirit of the Constitution, are Constitutional.''
\end{quote}
\noindent --- Chief Justice Marshall, Opinion of the Court, 17 U.S. (4 Wheat.) at 421.

\end{document}
